\section*{अध्याय \ref{ch:b2:atomic-orbitals} --- परमाण्वीय कक्षक: आकृतियाँ, ऊर्जाएँ और आकार}

\begin{solution}{exo:b2:atomic-orbitals:1}
3p: $n = 3$, $l = 1$; एक \omterm{def:b2:atomic-orbitals:node}{त्रिज्य नोड} ($n - l - 1$) और एक \omterm{def:b2:atomic-orbitals:node}{कोणीय नोड}। 4d: $n = 4$,
$l = 2$; एक \omterm{def:b2:atomic-orbitals:node}{त्रिज्य नोड} और दो \omterm{def:b2:atomic-orbitals:node}{कोणीय नोड}।
\end{solution}

\begin{solution}{exo:b2:atomic-orbitals:2}
$P_{2p} \propto r^4\mathrm e^{-r}$ ($r$, $a_0$ में): $d\ln P/dr = 4/r - 1 = 0$, $r = 4$ के लिए:
$4a_0 = \qty{212}{pm}$।
\end{solution}

\begin{solution}{exo:b2:atomic-orbitals:3}
$1 - \mathrm e^{-2}(1 + 2 + 2) = 1 - 5\mathrm e^{-2} = 0.323$।
\end{solution}

\begin{solution}{exo:b2:atomic-orbitals:4}
$x$ और $y$ अक्षों के बीच इंगित करने वाले चार पालि, $xy$ तल में, एकांतर
चिह्नों के (धनात्मक जहाँ $xy > 0$); नोडीय तल $xz$ और $yz$ हैं। \omterm{def:b2:atomic-orbitals:node}{त्रिज्य नोड}:
$n - l - 1 = 0$।
\end{solution}

\begin{solution}{exo:b2:atomic-orbitals:5}
ऑक्सीजन, $(1\mathrm s)^2(2\mathrm s, 2\mathrm p)^6$: $\sigma = 5 \times 0.35 + 2 \times 0.85
= 3.45$, $Z^* = 4.55$; त्रिज्या $4/4.55 = 0.88a_0$, \qty{47}{pm}।
\end{solution}

\begin{solution}{exo:b2:atomic-orbitals:6}
आयनन सबसे कम बँधे इलेक्ट्रॉनों को हटाता है। 4s इलेक्ट्रॉन
(\qty{-14}{eV}) 3d इलेक्ट्रॉनों (\qty{-59}{eV}) से कहीं कम बँधे हैं: दोनों 4s
इलेक्ट्रॉन पहले निकलते हैं, फिर एक 3d इलेक्ट्रॉन, जिससे $[\ce{Ar}]\,3\mathrm d^5$ मिलता है।
\end{solution}

\begin{solution}{exo:b2:atomic-orbitals:7}
लिथियम, 2s: $Z^* = 3 - 2 \times 0.85 = 1.30$, $I = 13.6 \times 1.30^2/4 =
\qty{5.75}{eV}$ (मापा गया 5.39, 7~\% बहुत ऊँचा)। सोडियम, 3s: $Z^* = 11 - 8 \times 0.85 -
2 \times 1.00 = 2.20$, $I = 13.6 \times 2.20^2/9 = \qty{7.3}{eV}$ (मापा गया 5.14):
नियम, जो एक स्थूल मॉडल हैं, बाहरी इलेक्ट्रॉन के बंधन को बढ़ा-चढ़ाकर आँकते हैं, समूह में
जितना नीचे उतना अधिक।
\end{solution}

\begin{solution}{exo:b2:atomic-orbitals:8}
सोडियम: $Z^* = 2.20$, त्रिज्या $9/2.20 = 4.1a_0$ (\qty{216}{pm})। क्लोरीन, 3p:
$\sigma = 6 \times 0.35 + 8 \times 0.85 + 2 \times 1.00 = 10.90$, $Z^* = 6.10$, त्रिज्या
$9/6.10 = 1.5a_0$ (\qty{78}{pm})। कोश वही है, परंतु क्लोरीन का बाहरी इलेक्ट्रॉन
लगभग तिगुना आवेश अनुभव करता है, और भीतर खिंच जाता है।
\end{solution}

\begin{solution}{exo:b2:atomic-orbitals:9}
दोनों का \omterm{def:b2:atomic-orbitals:radial-angular}{कोणीय भाग} $1/\sqrt{4\pi}$ है, इसलिए समाकल घटकर यह रह जाता है:
\[
  \int_0^\infty R_{1s}R_{2s}r^2\,dr \propto \int_0^\infty (2 - r)r^2\mathrm e^{-3r/2}\,dr =
  2 \times \frac{2!}{(3/2)^3} - \frac{3!}{(3/2)^4} = \frac{32}{27} - \frac{32}{27} = 0 .
\]
\end{solution}

\begin{solution}{exo:b2:atomic-orbitals:10}
$\int_0^\infty N^2r^2\mathrm e^{-r}r^2\,dr = N^2 \times 4! = 1$: $N = 1/\sqrt{24} =
1/(2\sqrt6)$, जैसा सारणी में है।
\end{solution}

\begin{solution}{exo:b2:atomic-orbitals:11}
$\langle r\rangle = 1.5a_0$, \omterm{def:b2:atomic-orbitals:radial-distribution}{सर्वाधिक संभाव्य त्रिज्या} $a_0$। $P(r) = 4r^2\mathrm e^{-2r}$ शून्य से
तीव्रता से उठता है और धीरे-धीरे गिरता है: बड़े $r$ पर लंबी पूँछ अधिकतम की
स्थिति की तुलना में माध्य में अधिक भार रखती है।
\end{solution}

\begin{solution}{exo:b2:atomic-orbitals:12}
$Y_{p_z} \propto z/r = \cos\theta$, $\theta = \pi/2$ के लिए शून्य होता है: $xy$ तल।
$\psi$ का चिह्न $z$ का है, ऊपर धनात्मक, नीचे ऋणात्मक। तल पर
\omterm{def:b2:atomic-orbitals:wavefunction}{प्रायिकता घनत्व} शून्य है।
\end{solution}

\begin{solution}{pb:b2:atomic-orbitals:1}
\textbf{1.} $\int|\psi|^2\,dV = \frac{4\pi}{\pi a_0^3}\int_0^\infty r^2\mathrm e^{-2r/a_0}\,dr =
\frac{4}{a_0^3} \times \frac{a_0^3}{4} = 1$।
\textbf{2.} $P(r) = 4\pi r^2|\psi|^2 = (4/a_0^3)\,r^2\mathrm e^{-2r/a_0}$।
\textbf{3.} $dP/dr = 0$: $2/r - 2/a_0 = 0$, $r = a_0$।
\textbf{4.} $\langle r\rangle = (4/a_0^3)\int_0^\infty r^3\mathrm e^{-2r/a_0}\,dr = (4/a_0^3)
\times 6(a_0/2)^4 = 1.5a_0$।
\textbf{5.} वितरण विषम है, बड़े $r$ पर लंबी पूँछ के साथ।
\textbf{6.} $|\psi|^2$ प्रति इकाई आयतन घनत्व है, जो नाभिक पर अधिकतम है; $P(r)$
कोश का आयतन, $4\pi r^2\,dr$, भी गिनता है, जो $r = 0$ पर शून्य होता है।
\textbf{7.} \qty{52.9}{pm}।
\textbf{8.} $x = r/a_0$ के साथ: $\int_{x_0}^\infty 4x^2\mathrm e^{-2x}\,dx$; दो बार खंडशः,
$= \mathrm e^{-2x_0}(2x_0^2 + 2x_0 + 1)$।
\textbf{9.} $5\mathrm e^{-2} = 0.677$।
\textbf{10.} $13\mathrm e^{-4} = 0.238$।
\textbf{11.} $\mathrm e^{-2x}(1 + 2x + 2x^2) = 0.10$, $x \approx 2.66$ के लिए: \qty{141}{pm}।
\textbf{12.} घनत्व किसी भी दूरी पर कभी शून्य नहीं होता: कोई किनारा नहीं। परंतु
प्रायिकता का 90~\% \qty{141}{pm} के भीतर है: एक व्यावहारिक आकार।
\textbf{13.} 2s: एक \omterm{def:b2:atomic-orbitals:node}{त्रिज्य नोड} ($2a_0$ त्रिज्या का गोला), कोई \omterm{def:b2:atomic-orbitals:node}{कोणीय नोड} नहीं। 2p: एक
\omterm{def:b2:atomic-orbitals:node}{कोणीय नोड} (एक तल), कोई \omterm{def:b2:atomic-orbitals:node}{त्रिज्य नोड} नहीं।
\textbf{14.} $r = 2a_0$ पर।
\textbf{15.} $4a_0$ (अभ्यास 2)।
\textbf{16.} $d\ln P/dr = 2/r - 2/(2 - r) - 1 = 0$ से $r^2 - 6r + 4 = 0$, $r = 3 \pm \sqrt5$:
$0.76a_0$ (छोटा भीतरी अधिकतम) और $5.24a_0$ (मुख्य अधिकतम)।
\textbf{17.} 2s, अपने भीतरी अधिकतम के कारण। बहु-इलेक्ट्रॉन परमाणु में 2s इलेक्ट्रॉन
1s कोश का भेदन करता है, 2p इलेक्ट्रॉन से कम आवरित होता है और ऊर्जा में
नीचे होता है।
\textbf{18.} सभी त्रिज्याएँ $Z = 6$ से भाग दी जाती हैं: 2p, $0.67a_0$ (\qty{35}{pm}) पर; 2s के अधिकतम
$0.13a_0$ और $0.87a_0$ पर।
\textbf{19.} C 3.25, N 3.90, O 4.55।
\textbf{20.} $\varepsilon = -13.6\,Z^{*2}/4$: C \qty{-35.9}{eV}, N \qty{-51.7}{eV}, O
\qty{-70.4}{eV}।
\textbf{21.} अध्याय की तरह: \qty{11.5}{eV} (मापा गया 11.26)।
\textbf{22.} N: परमाणु $5 \times (-13.6 \times 3.90^2/4) = \qty{-258.6}{eV}$; \ce{N+}, $Z^* =
4.25$, $4 \times (-13.6 \times 4.25^2/4) = \qty{-245.7}{eV}$; $I = \qty{12.9}{eV}$।
\textbf{23.} O: परमाणु $6 \times (-13.6 \times 4.55^2/4) = \qty{-422.3}{eV}$; \ce{O+}, $Z^* =
4.90$, $5 \times (-13.6 \times 4.90^2/4) = \qty{-408.2}{eV}$; $I = \qty{14.2}{eV}$। मॉडल
11.5, 12.9, 14.2; मापे गए 11.26, 14.53, 13.62।
\textbf{24.} ऑक्सीजन की मापी गई ऊर्जा नाइट्रोजन से कम है: ऑक्सीजन में
चौथे 2p इलेक्ट्रॉन को किसी दूसरे के साथ एक कक्षक साझा करना पड़ता है, और उनका
प्रतिकर्षण उसे हटाना आसान बनाता है, जबकि नाइट्रोजन के तीन अयुग्मित p इलेक्ट्रॉन
एक स्थायी अर्ध-पूरित उपकोश बनाते हैं। \omterm{def:b2:atomic-orbitals:slater}{स्लेटर के नियम} किसी समूह के सभी इलेक्ट्रॉनों को
समान मानते हैं और युग्मन के बारे में कुछ नहीं जानते।
\textbf{25.} हाइड्रोजन का 1s इलेक्ट्रॉन $2a_0$ से अधिक दूर
$\boldsymbol{13\mathrm e^{-4} \approx 0.238}$ प्रायिकता के साथ होता है।
\end{solution}
